\SourceLeaf{480}

\Needspace{5\baselineskip}\begin{center}APPENDIX.\end{center}
\BookRule{0.15\linewidth}

\Needspace{5\baselineskip}\begin{center}\small MANUSCRIPT NOTES BY GAUSS.\end{center}


On Article 40. \textit{With a third number \(C\) added, let \(\lambda'\) be the greatest common divisor of the numbers \(\lambda,C\), and determine numbers \(k,\gamma\) such that \(k\lambda+\gamma C=\lambda'\), whence \(k\alpha A+k\beta B+\gamma C=\lambda'\). It is clear that \(\lambda'\) is a common divisor of the numbers \(A,B.C\), and indeed the greatest; for if there were a greater one \(=\theta\), then \[k\alpha.\frac{A}{\theta}+k\beta.\frac{B}{\theta}+\gamma.\frac{C}{\theta}=\frac{\lambda'}{\theta}\quad\textit{integer, Q. E. A.}\]}


\textit{Thus what was proposed has been accomplished, on setting \(k\alpha=a,\ k\beta=b,\ \gamma=c,\ \lambda'=\mu\).}


On Article 114. \textit{The proof can be presented more elegantly as follows: \((a^{3n}-a^n)^2=2+(a^{4n}+1)(a^{2n}-2)\), \((a^{3n}+a^n)^2=-2+(a^{4n}+1)(a^{2n}+2)\), and hence \(\surd2\equiv\pm(a^{3n}-a^n),\ \surd-2\equiv\pm(a^{3n}+a^n)\ (\mathrm{mod.}\ 8n+1)\)}


In Article 256 VI the 16 positive determinants of the form \(8n+5\) for which the number of properly primitive classes is three times as large as that of improperly primitive classes, ‘37, .. 573’, are supplemented by: \textit{677, 701, 709, 757, 781, 813, 829, 877, 885, 901, 909, 925, 933, 973, 997, so that among 125 there are 31.}


On the words of Article 301: By this method the sum of the numbers of genera for determinants \(-\): up to \(-100\) is found \(=234,4\), whereas it is actually 233. — — \textit{From \(-1\) to \(-3000\), table 11166, formula 11167,9.}


On Article 336. \textit{If all numbers of the form \(2^{2^m}+1\) were prime, a sufficiently approximate expression for the number of the numbers in question \((N)\) smaller than a given number \(M\) would be \(\frac12\left(\frac{\log M}{\log2}\right)^2\).}


On the theorem in Article 42 on the divisors of a polynomial with integer coefficients — — \textit{1797 July 22.}


On the words of Article 130: After we have rigorously proved that every prime of the form \(4n+1\), taken both positively and negatively, is a nonresidue of some smaller prime, — — \textit{We discovered this proof in 1796, on April 8.}


On Article 131. \textit{The fundamental theorem was discovered by induction in March 1795. The first proof, given in this section, was found in April 1796.}

\SourceLeaf{481}


On the words of Article 133: Let us undertake the investigation (of the fundamental theorem) still more generally. Consider any two odd coprime numbers, with any signs, \(P\) and \(Q\). — — \textit{1796 April 29.}


On the words of Article 145: Moreover, the theorems concerning the residues \(+2\) and \(-2\) would then have had to be assumed; but since our proof has been completed without these theorems, we obtain from it a new method of proving them. — — \textit{1797 February 4.}


On the heading of Section V: Forms and indeterminate equations of the second degree. — — \textit{From June 22, 1796 onward.}


On the words of Article 234: .. we pass to another subject of the greatest importance, hitherto untouched by anyone, the composition of forms. — — \textit{These investigations were begun in the autumn of 1798.}


On the words of Article 262: From this principle we can derive a new method of proving not only the fundamental theorem, but also the other theorems of the preceding section concerning the residues \(-1,\allowbreak +2,\allowbreak -2\), — — \textit{The principles of this method first presented themselves in 1796, on July 27, but were developed and brought into their present form in the spring of 1800.}


On the words of Article 266: .. but since several truths concerning them, and very beautiful ones, still remain, whose proper source must be sought in the theory of ternary forms of the second degree, we insert here a short digression on this theory, — — \textit{February 14, 1799.}


On the words of Article 272: namely, by showing first how any ternary form can be reduced to a simpler form, and then that the number of the simplest forms (reached by such reductions) is finite for any given determinant. — — \textit{1800 February 13.}


On the words of Article 287 III: In an entirely similar way it is proved that in the improperly primitive order all the characters found to be the only possible ones by the rules of Article 264 II, III are possible, whether they are \(P\) or \(Q\). — These theorems, etc. — — \textit{They were first supplied with proofs in April 1798.}


On the words of Article 302: The average number of classes, however (which will need no definition), grows very regularly, — — \textit{The first idea at the beginning of 1799.}


On the words of Article 306 X: Finally, we observe that, since all the properties considered in this and the preceding article depend chiefly on the number \(n\), which is something similar to \(p-1\) in Section III, this number deserves the closest attention; it would therefore be most desirable to discover a general relation between it and the determinant to which it belongs. — — \textit{It has turned out as we wished—and so completely that nothing further remains to be desired. November 30—December 3, 1800.}


On Article 365. \textit{We discovered that the circle can be divided geometrically into 17 parts in 1796, on March 30.}

\BookRule{0.3\linewidth}

\clearpage
